\begin{abstract}
Auto-regressive token generation in large language models is memory-bound because it requires ``attending to'' key and value tensors (KV cache) of all previous tokens. Prior work aims to improve the efficiency of this \textit{decode} process by batching multiple requests together, and maximizing batch size subject to GPU memory constraints. The key observation of our work is that with prefix-sharing workloads, smaller, prefix-homogeneous batches -- where \textit{all} requests share a common prefix -- can achieve higher decode throughput than larger, heterogeneous batches, due to better spatial and temporal locality during KV cache accesses. However, prefix-aware schedulers in state-of-the-art inference engines maximize prefix reuse within a batch only to reduce KV cache memory footprint, but do not stop batch formation at smaller homogeneous batches that could have performed better. Further, we show that shared prefix detection in existing schedulers relies on radix-tree traversals, incurring substantial CPU overhead that is often comparable to GPU execution time. This paper presents \Feather, a prefix-aware scheduler that uses reinforcement learning (RL) to learn the optimal tradeoff between batch size and prefix homogeneity. We also introduce Chunked Hash Tree (CHT), a lightweight data structure that enables fast prefix detection and efficient request selection for the RL scheduler, avoiding expensive tree traversals. We integrate \Feather into vLLM and SGLang, and our evaluation shows that \Feather achieves 2--10$\times$ higher end-to-end throughput as compared to existing schedulers, while doing no worse than the status quo when the workload does not have enough prefix sharing. \Feather achieves these gains by reducing the total number of KV cache accesses, surpassing the performance of prefix-aware attention kernels that have the same goal.
\end{abstract}